% Compile with pdflatex (or tectonic); uses the official ICML 2026 style files.

\documentclass{article}

\usepackage{microtype}
\usepackage{graphicx}
\usepackage{subcaption}
\usepackage{booktabs}

\usepackage{hyperref}

% Preprint (non-blind) build.
\usepackage[preprint]{icml2026}

\usepackage{amsmath}
\usepackage{amssymb}
\usepackage{mathtools}

% Bra-ket notation, defined locally so we don't depend on the braket package.
\newcommand{\ket}[1]{\left|#1\right\rangle}
\newcommand{\bra}[1]{\left\langle#1\right|}
\newcommand{\braket}[1]{\left\langle#1\right\rangle}

\usepackage[capitalize,noabbrev]{cleveref}

\icmltitlerunning{Can a Quantum Circuit Read a Sentence?}

\begin{document}

\twocolumn[
  \icmltitle{Can a Quantum Circuit Read a Sentence?}

  \icmlsetsymbol{equal}{*}

  \begin{icmlauthorlist}
    \icmlauthor{Feolu Kolawole}{equal,stan}
    \icmlauthor{Nick Rui}{equal,stan}
  \end{icmlauthorlist}

  \icmlaffiliation{stan}{Physics 14N, Stanford University, Stanford, CA, USA}

  \icmlcorrespondingauthor{Feolu Kolawole}{flukol@stanford.edu}
  \icmlcorrespondingauthor{Nick Rui}{nickrui@stanford.edu}

  \icmlkeywords{Quantum Natural Language Processing, DisCoCat, Variational Quantum Circuits}

  \vskip 0.3in
]

\printAffiliationsAndNotice{\icmlEqualContribution}

% Force a clean running header (the style suppresses any single \small line
% taller than 6.25pt, which Computer Modern's cap height exceeds).
\makeatletter\gdef\@icmltitlerunning{Can a Quantum Circuit Read a Sentence?}\makeatother

\begin{abstract}
  Quantum natural language processing (QNLP) rests on a mathematical
  coincidence. Grammar composes the meaning of a sentence in the same way that
  qubits compose into a larger quantum state. A sentence can therefore be read
  as a recipe for a quantum circuit, and the circuit's measurement can be read
  as the sentence's meaning. We study this idea at the smallest scale where
  every step can be checked by hand. Our model is a sentiment classifier that
  maps a short sentence to a two-qubit circuit. Grammar decides which word
  rotates which qubit, an entangling gate ties the subject to the predicate, and
  one measurement returns positive or negative. Each prediction is the result of
  simulating the circuit exactly. The parameters are trained with the
  parameter-shift rule, the analytic-gradient method also used on quantum
  hardware. On a small balanced dataset the classifier labels every sentence
  correctly. We trace one example through the full pipeline, from parse to final
  amplitude. The goal is understanding rather than performance. The paper gives a
  working demonstration of how syntax can drive a quantum circuit, and a clear
  account of what two qubits can and cannot show.
\end{abstract}

\section{Introduction}
\label{sec:intro}

Most natural language processing learns vector representations of words from
large corpora. Grammar, if it appears at all, is something the model discovers
on its own. QNLP starts from the other side. In the DisCoCat model of
\citet{coecke2010mathematical}, short for distributional compositional
categorical, grammar comes first. A noun is a state. A verb is a process that
consumes the nouns around it. The grammatical reduction of a sentence specifies
how to combine those pieces. Read literally, that diagram is a quantum circuit
\citep{coecke2020foundations}. The analogy is more than a slogan because
compositional semantics and multi-qubit Hilbert space are both built from tensor
products. The bookkeeping is the same on both sides.

We wanted to see this work in practice. We built the smallest model that still
does something useful: a binary sentiment classifier for short sentences such as
``the clever program runs.'' It uses two qubits, a closed vocabulary, and three
sentence templates. At this size the circuit can be simulated exactly, trained
the way a real QNLP circuit would be, and checked line by line.

Two details keep the exercise honest. First, the prediction is read off a
simulated quantum state through the Born rule, not off a classical sigmoid placed
next to a circuit diagram. Second, the parameters are trained with the
parameter-shift rule \citep{mitarai2018quantum,schuld2019evaluating}, the exact
gradient method required on physical hardware. We are also clear about scale. Two
qubits can show compositionality at work. They cannot demonstrate a quantum
advantage. \Cref{sec:limits} states where the hard problems begin.

\section{Background}
\label{sec:background}

\subsection{Qubits and Measurement}

A classical bit is either $0$ or $1$. A qubit is a unit vector in $\mathbb{C}^2$,
written in the computational basis as
\begin{equation}
  \ket{\psi} = \alpha\ket{0} + \beta\ket{1}, \qquad |\alpha|^2 + |\beta|^2 = 1.
\end{equation}
The only single-qubit gate we need is the $R_y$ rotation,
\begin{equation}
  R_y(\theta) =
  \begin{pmatrix} \cos\tfrac{\theta}{2} & -\sin\tfrac{\theta}{2} \\[2pt]
                  \sin\tfrac{\theta}{2} & \phantom{-}\cos\tfrac{\theta}{2} \end{pmatrix},
  \label{eq:ry}
\end{equation}
which sends $\ket{0}$ to $\cos(\theta/2)\ket{0} + \sin(\theta/2)\ket{1}$. This is
a point at polar angle $\theta$ below the north pole of the Bloch sphere.
Measuring in the $\{\ket{0},\ket{1}\}$ basis collapses the state onto one pole.
The probabilities follow the Born rule,
\begin{equation}
  P(\ket{1}) = |\!\braket{1|\psi}\!|^2 = \sin^2(\theta/2),
  \label{eq:born}
\end{equation}
so a state on the equator ($\theta = \pi/2$) is a fair coin. We use only real
amplitudes and $R_y$ gates. The whole computation then lives on a single great
circle of the Bloch sphere, which is the right level of machinery for a first
look at QNLP \citep{nielsen2010quantum}.

\subsection{Two Qubits and Entanglement}

Two qubits live in
$\mathcal{H}_\mathrm{subj} \otimes \mathcal{H}_\mathrm{sent} \cong \mathbb{C}^4$,
with basis $\{\ket{00},\ket{01},\ket{10},\ket{11}\}$. The left digit is the
subject qubit and the right digit the sentence qubit, so amplitude $i$
corresponds to $i = 2\,(\text{subject}) + (\text{sentence})$. A single-qubit gate
acts through a Kronecker product, for example $R_y(\theta)\otimes I$ on the
subject. The controlled-NOT gate, with the subject as control, flips the sentence
qubit when the subject is $\ket{1}$. It swaps $\ket{10}\leftrightarrow\ket{11}$:
\begin{equation}
  \mathrm{CX} =
  \begin{pmatrix}
    1&0&0&0\\ 0&1&0&0\\ 0&0&0&1\\ 0&0&1&0
  \end{pmatrix}.
  \label{eq:cx}
\end{equation}
When both qubits are already in superposition, this gate produces a state that
does not factor as $\ket{\psi_\mathrm{subj}}\otimes\ket{\psi_\mathrm{sent}}$. The
two qubits become entangled. This is the behavior we want from grammatical
composition, where subject and predicate can no longer be read in isolation.

\subsection{Grammar in DisCoCat}

DisCoCat assigns each word a grammatical type built from a noun type $n$ and a
sentence type $s$ \citep{coecke2010mathematical}. A transitive verb carries type
$n^r\, s\, n^l$. It expects a noun on each side and returns a sentence. Reducing
the sentence contracts the matching wires, called cups, and the surviving wire is
the sentence meaning. Reading the diagram as a quantum process turns nouns into
states, verbs into multi-linear maps, and cups into Bell-style contractions
\citep{coecke2020foundations}. Our model keeps the central idea, that grammar
fixes how meanings combine before any number is computed. It reduces the
apparatus to two qubits so the whole thing can be simulated and trained exactly.

\section{The Model}
\label{sec:model}

\subsection{Constrained Grammar}

To remove any ambiguity, we restrict every sentence to one of three templates
over a fixed vocabulary:
\begin{center}\small
\begin{tabular}{ll}
  \texttt{the ADJ NOUN VERB} & intransitive \\
  \texttt{the ADJ NOUN VERB the NOUN} & transitive \\
  \texttt{the ADJ NOUN VERB the ADJ NOUN} & + obj.\ adj.
\end{tabular}
\end{center}
A small parser reads each sentence into its grammatical roles: a subject noun, a
subject adjective, a verb, and optionally an object noun and its adjective. This
replaces the CCG and DisCoCat machinery that a library such as lambeq
\citep{kartsaklis2021lambeq} would provide at scale. We give up coverage for a
parse we can trust and reproduce. \Cref{fig:grammar} shows one parse.

\subsection{Encoding Words as Angles}

Each content word $w$ carries one trainable number $\phi_w$, initialized near
zero. Grammar enters twice. First, it decides which qubit a word turns. The
subject-side words, the adjective and subject noun, sum into the subject qubit's
rotation. The predicate-side words, the verb and any object or object adjective,
sum into the sentence qubit's rotation:
\begin{align}
  \theta_\mathrm{subj} &= \phi_\mathrm{adj} + \phi_\mathrm{subj}, \\
  \theta_\mathrm{pred} &= \phi_\mathrm{verb} + \phi_\mathrm{obj\text{-}adj} + \phi_\mathrm{obj}.
\end{align}
Absent roles drop out of the sums.

Second, grammar decides how strongly each role contributes. A role-weighted sum
forms a single sentiment score. Verbs and adjectives carry most of the sentiment
in this domain, so they receive larger weights than the nouns:
\begin{equation}
  s = b + \sum_{(\text{role},\,w)\in\text{roles}} \omega_\mathrm{role}\,\phi_w,
  \label{eq:score}
\end{equation}
where the role weights $\omega$ appear in \cref{tab:weights} and $b$ is a
trainable bias. The score sets the final rotation, offset by $\pi/2$ so that a
neutral sentence leaves the sentence qubit on the equator:
\begin{equation}
  \theta_\mathrm{dec} = \frac{\pi}{2} + s.
  \label{eq:decision}
\end{equation}

\begin{table}[t]
  \caption{Role weights $\omega_\mathrm{role}$ in the sentiment score of
  \cref{eq:score}. Verbs and subject adjectives carry the sentiment here and are
  weighted accordingly. Nouns contribute little. Object weights apply only when
  the sentence has an object.}
  \label{tab:weights}
  \vskip 0.1in
  \centering
  \begin{tabular}{lc}
    \toprule
    Role & Weight $\omega$ \\
    \midrule
    Verb              & 1.40 \\
    Subject adjective & 1.25 \\
    Object adjective  & 0.90 \\
    Subject noun      & 0.45 \\
    Object noun       & 0.35 \\
    \bottomrule
  \end{tabular}
  \vskip -0.1in
\end{table}

\subsection{Circuit and Readout}

The circuit (\cref{fig:circuit}) runs four steps. It applies
$R_y(\theta_\mathrm{subj})$ to the subject qubit and $R_y(\theta_\mathrm{pred})$
to the sentence qubit to load the words. It applies a CX with the subject as
control to compose them. It applies $R_y(\theta_\mathrm{dec})$ to the sentence
qubit for the verdict. It measures the sentence qubit. Starting from $\ket{00}$,
the state after all four gates is
\begin{multline}
  \ket{\Psi} =
  (I\otimes R_y(\theta_\mathrm{dec}))\,\mathrm{CX} \\
  \times (R_y(\theta_\mathrm{subj})\otimes R_y(\theta_\mathrm{pred}))\,
  \ket{00}.
  \label{eq:forward}
\end{multline}
The model's only output is the probability of reading $1$ on the sentence qubit,
\begin{equation}
  P(\ket{1}) = |\Psi_{01}|^2 + |\Psi_{11}|^2.
  \label{eq:readout}
\end{equation}
We call $\ket{1}$ positive and $\ket{0}$ negative, and classify by whether
$P(\ket{1})$ exceeds one half. No separate scoring function runs offstage. The
circuit we draw and the prediction we make are the same code.

\begin{figure}[t]
  \centering
  \includegraphics[width=\columnwidth]{figures/grammar_diagram.png}
  \caption{The grammatical parse of an example sentence. These roles, not the
  order the words happen to arrive in, decide how the angles are wired into the
  circuit of \cref{fig:circuit}.}
  \label{fig:grammar}
\end{figure}

\begin{figure}[t]
  \centering
  \includegraphics[width=\columnwidth]{figures/toy_circuit.png}
  \caption{The two-qubit circuit. Word parameters set the opening rotations. The
  CX gate entangles subject and sentence, which is grammatical composition. A
  final $R_y(\theta_\mathrm{dec})$ writes in the sentiment verdict before the
  sentence qubit is measured.}
  \label{fig:circuit}
\end{figure}

\subsection{Training}
\label{sec:training}

We train by minimizing binary cross-entropy between $P(\ket{1})$ and the label.
The gradients are the interesting part. Each angle enters through a gate of the
form $e^{-i\theta G/2}$. For such gates the derivative of a measured probability
is given exactly by the parameter-shift rule
\citep{mitarai2018quantum,schuld2019evaluating}:
\begin{equation}
  \frac{\partial P}{\partial\theta}
  = \tfrac12\!\left[\,
      P\!\left(\theta + \tfrac{\pi}{2}\right)
    - P\!\left(\theta - \tfrac{\pi}{2}\right)
    \right].
  \label{eq:psr}
\end{equation}
This is not a finite-difference approximation. It is exact for these gates, and
it is how real QNLP circuits are differentiated on hardware. We evaluate it for
each of the three gate angles by re-running the simulation at shifted values. The
chain rule then carries those gradients back to the word parameters and the bias.
One word may feed several gates, weighted by how often it appears and by its role.
Gradient descent does the rest, over fifty epochs at a learning rate of $0.35$.

\section{Worked Example}
\label{sec:walkthrough}

To see the model run, take the positive sentence ``the clever program runs.''
After training, its words have settled on $\phi_\text{clever}=0.098$,
$\phi_\text{program}=0.043$, and $\phi_\text{runs}=0.489$, with bias $b=0.068$.

The subject qubit collects the adjective and noun,
$\theta_\mathrm{subj}=0.098+0.043=0.141$. The sentence qubit collects the verb,
$\theta_\mathrm{pred}=0.489$. The role-weighted score from \cref{eq:score} is
\begin{align*}
  s &= 0.068 + 1.25(0.098) + 0.45(0.043) + 1.40(0.489) \\
    &= 0.894,
\end{align*}
so the verdict rotation is $\theta_\mathrm{dec}=\pi/2 + 0.894 = 2.465$. The verb,
at weight $1.40$, dominates the score, which is the bias we built in.

Following the four amplitudes of \cref{eq:forward} through the circuit gives
\cref{tab:evolution}. The opening rotations barely disturb $\ket{00}$. The CX
swaps the $\ket{10}$ and $\ket{11}$ amplitudes, which is grammatical composition
in action. The verdict rotation tips the sentence qubit across the equator, from
the $\ket{0}$ hemisphere into the $\ket{1}$ hemisphere. Reading off
\cref{eq:readout},
\begin{equation*}
  P(\ket{1}) = 0.993^2 + 0.039^2 = 0.988
  \;\Rightarrow\; \textbf{POSITIVE}.
\end{equation*}
The opposite sentence, ``the broken code fails,'' runs through the same circuit
with its own angles and gives $P(\ket{1}) = 0.003$, a confident negative. Almost
the entire gap between the two comes from the verb and adjective angles, as the
role weights intend.

\begin{table}[t]
  \caption{The statevector for ``the clever program runs'' at each stage of
  \cref{eq:forward}. The left digit is the subject qubit, the right digit the
  sentence qubit. The $\ket{01}$ and $\ket{11}$ columns are the ones a positive
  reading draws on.}
  \label{tab:evolution}
  \vskip 0.1in
  \centering
  \small
  \setlength{\tabcolsep}{4pt}
  \begin{tabular}{lrrrr}
    \toprule
    Operation & $\ket{00}$ & $\ket{01}$ & $\ket{10}$ & $\ket{11}$ \\
    \midrule
    init $\ket{00}$            & 1.000 & 0.000 & 0.000 & 0.000 \\
    word $R_y$ (both qubits)   & 0.968 & 0.241 & 0.068 & 0.017 \\
    CX: $\ket{10}\!\leftrightarrow\!\ket{11}$ & 0.968 & 0.241 & 0.017 & 0.068 \\
    decision $R_y(2.465)$      & 0.094 & 0.993 & $-0.059$ & 0.039 \\
    \bottomrule
  \end{tabular}
  \vskip -0.1in
\end{table}

\section{Results}
\label{sec:experiments}

The dataset has 72 sentences, split evenly into 36 positive and 36 negative, with
60 for training and 12 for testing. Positive and negative sentences use identical
templates and differ only in their content words, such as \emph{clever} or
\emph{reliable} against \emph{broken} or \emph{buggy}. This sharpens the
question. With grammar fixing the wiring, can the learned word angles alone carry
the meaning? The splits, the parser, and the random seed are all deterministic,
so the run reproduces from a single command.

It can. The model labels every training and test sentence correctly
(\cref{tab:results}). The cross-entropy falls smoothly across the fifty epochs of
training (\cref{fig:loss}). Every prediction comes from exact statevector
simulation of the two qubits. At this size a forward pass and its gradients
finish in well under a second on a laptop.

Perfect accuracy here is not a strong claim. The dataset is small, grammatically
uniform, and separable in the feature the model was built around. A classical
linear classifier would also solve it. The experiment checks something narrower.
It checks that the quantum realization holds together. Rotations, entanglement, a
Born-rule readout, and parameter-shift training reproduce the correct meaning
from end to end, with grammar rather than word order doing the composing. The
result is a working mechanism, not a benchmark.

\begin{table}[t]
  \caption{Accuracy. Every prediction is the outcome of simulating and measuring
  the two-qubit circuit (\cref{eq:readout}).}
  \label{tab:results}
  \vskip 0.1in
  \centering
  \begin{tabular}{lcc}
    \toprule
    Split & \# Sentences & Accuracy \\
    \midrule
    Train & 60 & 100\% \\
    Test  & 12 & 100\% \\
    \bottomrule
  \end{tabular}
  \vskip -0.1in
\end{table}

\begin{figure}[t]
  \centering
  \includegraphics[width=\columnwidth]{figures/loss_curve.png}
  \caption{Binary cross-entropy over training. The gradients behind this curve
  come from the parameter-shift rule of \cref{eq:psr}, the same method used on a
  real quantum device.}
  \label{fig:loss}
\end{figure}

\section{Limitations}
\label{sec:limits}

The ceiling is worth stating plainly. Two qubits and a fixed grammar cannot
demonstrate a quantum advantage, and we do not claim one. Simulating a circuit
this small is trivial on a classical computer. The value is the demonstration
itself. A sentence's grammar can be compiled into a circuit whose measurement
encodes meaning, with no hidden steps.

Real QNLP changes the picture in ways this model does not capture. The
hand-written parser becomes a full DisCoCat or CCG pipeline
\citep{kartsaklis2021lambeq}. Words spread across several qubits each. The
circuits run on noisy devices where probabilities are sampled rather than
computed, so the analytic gradients of \cref{eq:psr} become a necessity rather
than a convenience \citep{lorenz2023qnlp}. New difficulties appear that do not
exist at two qubits. The sharpest is the barren plateau, where training gradients
shrink exponentially as qubits are added \citep{mcclean2018barren}. This is a
central obstacle for variational quantum models in general
\citep{cerezo2021variational}.

The modeling shortcuts also have a cost. Real amplitudes, one rotation per word,
a fixed table of role weights, and a closed vocabulary each trade expressivity
for clarity. Richer QNLP encodes words as multi-qubit states and verbs as genuine
multi-linear maps, so that composition is more than a single scalar score.
Extending the model in that direction, without losing the transparency that makes
it useful, is the natural next step.

\section{Conclusion}
\label{sec:conclusion}

We built a two-qubit sentiment classifier in which grammar lays out the circuit,
learned word angles supply the meaning, and one measurement gives the verdict.
Training uses the parameter-shift rule throughout. On a small balanced dataset
the model is perfectly accurate. We traced a single sentence through every stage,
from parse and angles to the entangling gate, the amplitudes, and the readout.
The model is small by design. It offers a clear picture of how syntax can drive a
quantum circuit, and a clear sense of how much further there is to go.

The dataset, parser, simulator, trainer, and figures are released with the
project, and the full pipeline runs deterministically from one command.

\bibliography{references}
\bibliographystyle{icml2026}

\end{document}
